\scanpage{181}\mspage{36}%
% archive sheet number: 181
% One continuous running paragraph; author number at top centre, archive pencil number at lower left.
pouvoir se paraphraser dans le cadre des schémas excellents, à grands coups de voisinages tubulaires étales.
Dans l'un et l'autre cas, le ``type d'homotopie fin'' d'un espace modéré, resp. d'un schéma excellent, se définirait par passage à la limite à partir des
% L7a: the opening quotation marks stand to the left of the raised prefix; its guide reaches the start of types.
``\inserted{$\unsure{\mathcal{P}|P}$-}types d'homotopie'' associés à des stratifications équisingulières (à strates non sing.) $\mathcal{P}$ de plus en plus fines.
Ce type d'homotopie \inserted{``fin''} incarnerait la connaissance, non des seuls faisceaux ou \inserted{$n$-}champs \unsure{loc.|loc} constants, mais
% L12a: complete interlinear parenthesis; the long guide descends before la.
\inserted{(via un passage à la limite \unsure{inductive})} la connai%
% L12 nontext point at approximately (714,1189), below the ai region; origin uncertain.
\raisebox{0pt}[0pt][0pt]{\parbox{0pt}{\figure{\begin{tikzpicture}[overlay,x=1cm,y=1cm]\fill (-0.04,-0.28) circle (0.35pt);\end{tikzpicture}}}}%
ssance de
% The visible rule begins within the o of tous, approximately (1171,1167), and ends near (1289,1164).
t\uline{ous}. Et il dépendrait, en un sens convenable,
fonctoriellement%
% L14 nontext point at approximately (700,1366), below/right of fonctoriellement; origin uncertain.
\raisebox{0pt}[0pt][0pt]{\parbox{0pt}{\figure{\begin{tikzpicture}[overlay,x=1cm,y=1cm]\fill (0.16,-0.12) circle (0.3pt);\end{tikzpicture}}}}
de $X$. Dans%
% L14 nontext point at approximately (1170,1369), below the end of Dans; origin uncertain.
\raisebox{0pt}[0pt][0pt]{\parbox{0pt}{\figure{\begin{tikzpicture}[overlay,x=1cm,y=1cm]\fill (0,-0.27) circle (0.45pt);\end{tikzpicture}}}}
le cas d'un%
% L14 nontext point at approximately (1627,1352), below/right of d'un, near the right margin; origin uncertain.
\raisebox{0pt}[0pt][0pt]{\parbox{0pt}{\figure{\begin{tikzpicture}[overlay,x=1cm,y=1cm]\fill (0.45,-0.13) circle (0.25pt);\end{tikzpicture}}}}
schéma de type fini sur un corps alg. clos disons, les
% Literal ruled span: from within the é of théorèmes, approximately (671,1580), to (1328,1557), before cohomologique.
th\uline{éorèmes de finitude} cohomologique ou homotopique les plus forts s'exprimeraient précisément en termes des types d'homotopie fin, et diraient que les types d'homotopie \unsure{ordinaire|ordïnaire} qui en sont les élément constituants sont essentiellement des ``\unsure{polyédres|polyèdres} finis'' -- et $\overline{m}$ des variétés à bord compactes -- ou de façon plus précise,
